% Part: first-order-logic 
% Chapter: axiomatic-deduction 
% Section: rules-and-proofs

\documentclass[../../../include/open-logic-section]{subfiles}

\begin{document}

\iftag{FOL}
      {\olfileid[pt]{fol}{axd}{rul}}
      {\olfileid[pt]{pl}{axd}{rul}}

\olsection{Regras e \usetoken{P}{derivation}}

\begin{explain}
  As !!{derivation}s axiomáticas talvez formem o sistema de
  !!{derivation} mais simples para a lógica. !!^a{derivation} é apenas
  uma sequência de !!{formula}s. Para contar como !!a{derivation}, cada
  !!{formula} da sequência deve ser uma instância de um axioma ou decorrer,
  por uma regra de inferência, de uma ou mais !!{formula}s que a precedem
  na sequência. !!^a{derivation} !!{derive}s sua última !!{formula}.
\end{explain}

\begin{defn}[!!^{derivability}]
Se $\Gamma$ é um conjunto de !!{formula}s de $\Lang L$, então uma
\emph{!!{derivation}} a partir de~$\Gamma$ é uma sequência finita $!A_1$,
\dots,~$!A_n$ de !!{formula}s tal que, para cada $i \le n$, ocorre uma
das seguintes possibilidades:
\begin{enumerate}
\item $!A_i \in \Gamma$; or
\item $!A_i$ é um axioma; ou
\item $!A_i$ decorre, por uma regra de inferência, de algum $!A_j$ (e
  $!A_k$) com $j < i$ (e $k < i$).
\end{enumerate}
\end{defn}

O que conta como uma !!{derivation} correta depende das regras de
inferência que admitimos (e, naturalmente, do que tomamos como axiomas).
Uma regra de inferência é uma afirmação condicional que nos diz que, sob
certas condições, um passo~$A_i$ em !!a{derivation} é um passo de
inferência correto.

\begin{defn}[Regra de inferência]
Uma \emph{regra de inferência} fornece uma condição suficiente para que
um passo conte como passo de inferência correto em !!a{derivation} a
partir de~$\Gamma$.
\end{defn}

Por exemplo, como toda sequência de um elemento $!A$ com $!A \in \Gamma$
conta trivialmente como !!a{derivation}, a seguinte pode ser uma regra de
inferência muito simples:
\begin{quote}
  Se $!A \in \Gamma$, então $!A$ é sempre um passo de inferência correto
  em qualquer !!{derivation} a partir de~$\Gamma$.
\end{quote}
De modo semelhante, se $!A$ é um dos axiomas, então $!A$ por si só é
!!a{derivation}; portanto, esta também é uma regra de inferência:
\begin{quote}
  Se $!A$ é um axioma, então $!A$ é um passo de inferência correto.
\end{quote}
O caso fica mais interessante quando a regra de inferência recorre a
!!{formula}s que aparecem antes do passo considerado. A regra seguinte
se chama \emph{modus ponens:}
\begin{quote}
  Se $!B \lif !A$ e $!B$ ocorrem anteriormente na !!{derivation},
  então~$!A$ é um passo de inferência correto.
\end{quote}
Se esta for a única regra de inferência, então nossa definição de
!!{derivation} acima equivale ao seguinte: $!A_1$, \dots,~$!A_n$ é
!!a{derivation} se, e somente se, para cada $i \le n$ ocorre uma das
seguintes possibilidades:
\begin{enumerate}
\item $!A_i \in \Gamma$; or
\item $!A_i$ é um axioma; ou
\item para algum $j < i$, $!A_j$ é $!B \lif !A_i$ e, para algum $k < i$,
  $!A_k$ é~$!B$.
\end{enumerate}
A última cláusula diz que $!A_i$ decorre de~$!A_j$ ($!B \lif !A_i$) e
$!A_k$ ($!B$) por modus ponens. Se podemos avançar de $1$ até~$n$ e,
a cada passo, encontramos !!a{formula}~$!A_i$ que está em~$\Gamma$, é um
axioma ou é declarada passo de inferência correto por uma regra de
inferência, então a sequência inteira conta como !!{derivation} correta.

\begin{defn}[!!^{derivability}]
Uma !!{formula}~$!A$ é \emph{!!{derivable}} a partir de $\Gamma$, o que
escrevemos $\Gamma \Proves !A$, se existe !!a{derivation} a partir de
$\Gamma$ que termina em $!A$.
\end{defn}

\begin{defn}[Teoremas]
!!^a{formula}~$!A$ é um \emph{teorema} se existe !!a{derivation}
de~$!A$ a partir do conjunto vazio. Escrevemos $\Proves !A$ se $!A$ é
um teorema e $\Proves/ !A$ se não é.
\end{defn}


\end{document}
